% !TeX root = ../seminar.tex
\section{Preliminary Results}

\begin{frame}{What the first overview forecast covers}
\begin{columns}[T]
\begin{column}{.47\linewidth}
\textbf{Three simulated lens populations}
\begin{itemize}
  \item Stellar population with remnants
  \item Isolated black holes, $3$--$1000\,\Msun$
  \item Neutron stars, near $1.35\,\Msun$
\end{itemize}
\end{column}
\begin{column}{.49\linewidth}
\textbf{September 2026 runs}
\begin{itemize}
  \item $1{,}829$ sightlines over $68$\,deg$^2$
  \item Six circles for Roman, $1.47$\,deg$^2$
  \item A ten-year observing window
\end{itemize}
\end{column}
\end{columns}
\vspace{4mm}
\uncover<2->{%
\begin{beamercolorbox}[rounded=true,sep=6pt,wd=\linewidth]{block body}
\textbf{How the forecasts are reported}\par
We weight draws by event rate. Headline yields apply a dust adjustment after simulation.
For a compact-object abundance $F$, the event yield scales as $N(F)=F\,N_1$.
\end{beamercolorbox}}
\vfill
\refnote{First overview report and proposal. Yields use the per-resolved-object counting convention.}
\note{The ordinary-lens run uses a Kroupa initial mass function evolved for an old bulge. The separate black-hole
      run samples $dN/dM\propto M^{-1}$ and uses a simple kinematic model. The neutron-star distribution
      is centred on $1.35\,\Msun$ with width $0.15\,\Msun$. $F$ is the fraction of Galactic stellar mass
      assigned to the compact-object population. In September, extinction near the plane was diagnosed
      as too low. The headline yields reweight outcomes using expected losses at fainter magnitudes,
      with a reference dust map checked against the near-infrared VVV survey. The simulations themselves
      were left as run. The backups give the details and the external checks.}
\end{frame}

\begin{frame}{Roman detects more events within its footprint}
\begin{columns}[c]
\begin{column}{.48\linewidth}
\textbf{Ordinary lenses inside Roman's fields}\\[4pt]
\begin{tabular}{@{}lr@{}}
Roman detections & $\approx50{,}000$ \\
Rubin detections & $\approx4{,}900$
\end{tabular}
\vspace{6mm}

\uncover<2->{Outside those fields, Rubin accounts for \textbf{about 60\% of the detections across the simulated region}.}
\end{column}
\begin{column}{.47\linewidth}
\centering
\textbf{Extinction toward the bulge}\\[-1mm]
{\small At $A_V\approx7.2$ magnitudes}
\begin{tikzpicture}[font=\small]
  \fill[orange!65] (0,1.6) rectangle (5.18,2.1);
  \node[anchor=west, font=\footnotesize] at (0.1,1.85) {Rubin $r$: $6.1$\,mag};
  \fill[jdblue!65] (0,0.8) rectangle (1.19,1.3);
  \node[anchor=west, font=\footnotesize] at (1.3,1.05) {Roman F146: $1.4$\,mag};
  \draw (0,0.6) -- (0,2.25);
\end{tikzpicture}
\end{column}
\end{columns}
\vfill
\refnote{Rate-weighted, dust-adjusted yields for the ten-year window. Extinction law: \ads{1989ApJ...345..245C}{Cardelli et al.\ 1989}.}
\note{The two detection counts overlap. For ordinary lenses the report gives $50{,}350\pm2{,}945$
      for Roman and $4{,}853\pm515$ for Rubin within the footprint. Of all ordinary-lens detections
      in the simulated region, $(59\pm5)\%$ are outside Roman's footprint, where Rubin is the only
      simulated observer. The bar chart takes $A_V$ from the median toward Roman's footprint detections.
      Within the footprint, nearly all Rubin-only ordinary-lens detections peak when Roman is not observing.
      Rubin's time coverage matters most when the signal is concentrated inside a Roman gap.}
\end{frame}

\begin{frame}{Roman's data sharpen masses for events Rubin detects}
\centering
\textbf{Rubin detections inside Roman's fields}\\[3pt]
{\small Rubin data alone $\longrightarrow$ Rubin and Roman data combined}\\[3mm]
\renewcommand{\arraystretch}{1.2}
\begin{tabular}{lccc}
\toprule
 & ordinary lenses & neutron stars & black holes \\
\midrule
$\tE$ and $\piE$ above $2\sigma$ & $5\%\to28\%$ & $5\%\to37\%$ & $12\%\to83\%$ \\
\rowcolor{jdlightestblue}
mass error below $10\%$ & $\approx0\to8\%$ & $\approx0\to16\%$ & $0.2\%\to\mathbf{53\%}$ \\
\bottomrule
\end{tabular}

\vspace{4mm}
\uncover<2->{\small About \textbf{13\% more} black-hole masses reach $10\%$ in the joint forecast across the full footprint,
compared with the Roman-only yield.\par}
\vfill
\refnote{Rate-weighted and dust-adjusted. Underlying Rubin detections before dust adjustment have
$\Neff=685$ (ordinary), $1{,}687$ (neutron stars), and $2{,}011$ (black holes).}
\note{The table has a common denominator across each pair of fractions. The reverse comparison below it
      compares two total footprint yields, about $651F$ joint and $579F$ Roman-only for black holes.
      A $10\%$ mass means a forecast fractional $1\sigma$ error below $0.10$. The joint black-hole share
      in the table is $53.3\pm1.8$ percent. Most Rubin black-hole detections in these fields are also
      Roman detections. Rubin broadens both sky and time coverage, but extra visits often have little
      information on well-covered, heavily extinguished events.}
\end{frame}

\begin{frame}{For black holes, parallax limits the mass precision}
\begin{columns}[c]
\begin{column}{.52\linewidth}
\centering\small
\textbf{Roman precision in the September runs}
\begin{tikzpicture}[x=.046cm,y=.92cm,font=\footnotesize]
  \node[anchor=east] at (-3,2.3) {$\tetE$};
  \node[anchor=east] at (-3,1.3) {$\piE$};
  \node[anchor=east] at (-3,.3) {$\Ml$};
  \foreach \y/\val/\lab in {2.3/97/97,1.3/12.85/13,.3/12.6/13} {
    \fill[jdverylightgrey] (0,\y-.22) rectangle (100,\y+.22);
    \fill[jdblue!75] (0,\y-.22) rectangle (\val,\y+.22);
    \node[anchor=east] at (99,\y+.43) {$\lab\%$};
  }
\end{tikzpicture}\\[-1mm]
{\footnotesize Share with forecast error below $10\%$}
\end{column}
\begin{column}{.42\linewidth}
Roman's detections have a median maximum centroid shift of \textbf{$3.4$\,mas}.

\vspace{4mm}
\uncover<2->{The Earth--L2 baseline gives negligible median gain in black-hole parallax precision in the paired simulation.}
\end{column}
\end{columns}
\vfill
\refnote{Rate-weighted, as simulated. Each bar uses the Roman sample with a forecast available
for that parameter ($\Neff\simeq8{,}700$).}
\note{The rates shown here come from the report's compact-object table. The Roman-only mass
      fraction after the separate dust adjustment is about $10\%$ of Roman's detections.
      For the black-hole parallax test, rerun the precision calculation with Roman placed at Earth
      and keep the sampling and noise fixed. The median parallax-error ratio for L2 vs Earth was
      effectively one for black holes. The projected baseline is short, and the annual signal already
      contributes information across Roman's observing seasons.}
\end{frame}

\begin{frame}{How many black holes could Roman weigh?}
\centering
For a black-hole population containing \textbf{$0.5$--$3\%$} of Galactic stellar mass
\vspace{4mm}

\begin{columns}[T]
\begin{column}{.43\linewidth}
\begin{beamercolorbox}[rounded=true,sep=8pt,wd=\linewidth]{block body}
\centering{\Large\bfseries 30--170}\\[2pt]
black-hole events detected
\end{beamercolorbox}
\end{column}
\begin{column}{.43\linewidth}
\begin{beamercolorbox}[rounded=true,sep=8pt,wd=\linewidth]{block body example}
\centering{\Large\bfseries 3--17}\\[2pt]
masses with errors below $10\%$
\end{beamercolorbox}
\end{column}
\end{columns}
\vspace{4mm}

$\displaystyle N_{\rm det}=(5{,}660\pm320)F,\qquad N_{\rm mass}=(579\pm42)F$
\vspace{3mm}

\uncover<2->{\small Dust, field geometry, correlated astrometric errors and the assumed black-hole population all affect these yields.\par}
\vfill
\refnote{First overview report, after the dust adjustment. Roman-only forecast in the $1.47$\,deg$^2$
model footprint within the ten-year survey window.}
\note{The endpoints vary the assumed $F$ with the mass function and kinematics held fixed. They are an abundance
      scenario range. The uncertainties on the coefficients include Monte Carlo error and uncertainty
      in the dust adjustment. They do not include all survey and population-model uncertainties.
      An astrometric floor that is correlated across visits especially affects weaker centroid signals.
      The numerical yields here use the original model footprint, so the proposal includes a rerun on the adopted fields.}
\end{frame}
